\section{Effective Galois comparison: exact analytic inputs and constants}
\label{sec:effective-galois}

Throughout this section, $D_L=|\operatorname{disc}(L)|$, $n_L=[L:\mathbb Q]$, and all logarithms are natural.  Write $\mathrm{H}_M$ for the draft assertion that GRH holds for the Hecke $L$-functions of every number field.  Every conclusion using $\mathrm{H}_M$ remains conditional on that assertion.  The trivial Hecke character gives GRH for $\zeta_L$ for every number field $L$.

\paragraph{Source coverage.}
The six supplied \texttt{SORTED\_*.md} files are summaries, with frequently truncated statements.  The later recovery of the complete A1, successor, and LM8 papers permits the row-level scope to be checked against actual proofs.  In particular, LM8 A.32 uses Serre's qualitative open-image theorem only to obtain Zariski density of the commutator image.  Its actual factors are the verified non-CM models $F_{\pm6}/\mathbb Q$ of G06, with $j=-216$, subsequently restricted to the chosen number field $L$.  G06 supplies their conditional mod-$\ell$ bounds and unconditional adelic index bounds; no effective index was needed for the original qualitative conclusion.  LM8 D.25 uses Faltings's finiteness theorem for rational points, not the Faltings--Serre trace-comparison method.  The Gaussian and Fermat-cubic CM elliptic curves in other LM8 rows do not satisfy the non-CM hypotheses of G-LOMBARDO or G-MW1.  The row ledger records the precise finite fields in the recovered proofs.

\subsection{The finite-cover Chebotarev bound}

\paragraph{Cited theorem G-CHEB (Bach--Sorenson).}
If $L/K$ is finite Galois, $L\ne\mathbb Q$, and $C$ is a conjugacy class in $\operatorname{Gal}(L/K)$, then GRH for $\zeta_L$ and $\zeta_{\mathbb Q}$ gives an unramified degree-one prime $\mathfrak p$ of $K$ with Frobenius class $C$ and
\[
 N\mathfrak p\le B(L):=(4\log D_L+\tfrac52 n_L+5)^2.
\]
The exact reference is Bach--Sorenson, \emph{Explicit bounds for primes in residue classes}, Math. Comp. 65 (1996), 1717--1735, Theorem 5.1, p.~1729; the analytic hypotheses are specified in the last paragraph of \S2, p.~1720.  Their Theorem 3.2 and Corollary 3.3, p.~1721, make the reduction to a cyclic extension.  No Artin holomorphy conjecture is assumed.  This is already a conditional theorem of the literature; it is not a new finding of this work.

\paragraph{Finding G01 (a bounded prime outside a prescribed set).}
Let $L/\mathbb Q$ be finite Galois and let $S$ be a finite set of rational primes. Put
\[
 m=\operatorname{rad}\!\left(2\prod_{p\in S}p\right),\qquad
 \widetilde L=L(\sqrt m).
\]
For any nonempty conjugation-stable subset $C\subseteq\operatorname{Gal}(L/\mathbb Q)$, GRH for $\zeta_{\widetilde L}$ and $\zeta_{\mathbb Q}$ gives a rational prime $p\notin S\cup\{2\}$, unramified in $L$, such that
\[
 \operatorname{Frob}_p(L/\mathbb Q)\subseteq C,
 \qquad p\le B(\widetilde L).
\]
Thus $\mathrm{H}_M$ supplies this statement for arbitrary $L$. The stipulated Dirichlet GRH supplies it when $\widetilde L/\mathbb Q$ is abelian, in particular when $L/\mathbb Q$ is abelian.

\begin{proof}
The integer $m$ is even and squarefree, and $m\ge2$, so every prime dividing $m$ ramifies in the quadratic field $\mathbb Q(\sqrt m)$.  Therefore each such prime ramifies in $\widetilde L$.  Restriction $\operatorname{Gal}(\widetilde L/\mathbb Q)\to\operatorname{Gal}(L/\mathbb Q)$ is surjective.  The inverse image of $C$ is nonempty and stable under conjugation; choose one conjugacy class in it and apply G-CHEB.  The resulting prime is unramified in $\widetilde L$, hence does not divide $m$, and restriction of Frobenius proves the claim.  This argument is the squarefree-prime-avoidance construction of Mayle--Wang, Corollary~6, with $2$ explicitly included to remove a possible exceptional prime.
\end{proof}

\paragraph{Finding G02 (finite quotient trace comparison).}
Let $\rho_i:G_{\mathbb Q,S}\to\operatorname{GL}_d(R)$, $i=1,2$, be continuous representations over a finite commutative ring $R$, and let $L$ be the fixed field of the joint kernel. If their trace functions differ, they differ at a prime $p\notin S\cup\{2\}$ satisfying $p\le B(L(\sqrt m))$, under the GRH assumption in G01. The same assertion holds for any conjugation-invariant function of the joint image, including coefficients of characteristic polynomials.

\begin{proof}
The joint image is $\operatorname{Gal}(L/\mathbb Q)$.  The set on which the two trace functions differ is nonempty and stable under conjugation.  G01 applies to it.  A conjugation-invariant function is handled by precisely the same argument.
\end{proof}

\paragraph{Arbitrary number-field bases.}
For a finite Galois extension $L/K$, let $M/\mathbb Q$ be its normal closure and let $S_0$ contain the rational primes below every prime to be excluded.  With $m=\operatorname{rad}(2\prod_{p\in S_0}p)$, any prescribed class in $\operatorname{Gal}(L/K)$ occurs at a degree-one prime $\mathfrak p$ of $K$, above a rational prime outside $S_0\cup\{2\}$, with
\[
 N\mathfrak p\le B(M(\sqrt m))
\]
under GRH for $\zeta_{M(\sqrt m)}$ and $\zeta_{\mathbb Q}$.  To see this, extend a representative $g\in\operatorname{Gal}(L/K)$ to $\gamma\in\operatorname{Gal}(M/K)$ and apply G01 to its conjugacy class in $\operatorname{Gal}(M/\mathbb Q)$.  Choose a prime of $M$ whose Frobenius is precisely $\gamma$.  Since $\gamma$ fixes $K$, the prime below it in $K$ has residue degree one over $\mathbb Q$; its Frobenius in $L/K$ is $g$.  This gives an explicit form for the arbitrary-base finite affine images in the ledger, after their finite extension is fixed.  It does not bound a preceding Hilbert specialization or a depth parameter not specified by the proof.

This test compares trace functions at the chosen finite level.  Equality modulo a fixed power of a coefficient prime does not imply equality of the full adic representations.  To obtain a finite test for their full trace functions, one needs the next construction.

\subsection{A complete finite test for adic semisimplifications}

\paragraph{Finding G03 (an explicit uniform trace bound).}
Let $\mathcal O$ be the ring of integers in a finite extension $F/\mathbb Q_\ell$, with uniformizer $\varpi$ and residue field of cardinality $q$.  Let
\[
 \rho_1,\rho_2:G_{\mathbb Q,S}\longrightarrow\operatorname{GL}_d(\mathcal O)
\]
be continuous.  Set
\[
 N=2(q^{2d^2}-1),\qquad
 m=\operatorname{rad}\!\left(2\prod_{p\in S}p\right),
\]
and
\[
 \mathcal B(d,q,S)=
 \left[4\bigl((N-1)\log m+N\log N\bigr)+\tfrac52 N+5\right]^2.
 \tag{G03.1}
\]
Conditional on $\mathrm{H}_M$, equality of the exact Frobenius traces for every prime $p\nmid m$ with $p\le\mathcal B(d,q,S)$ is equivalent to equality of their trace functions on $G_{\mathbb Q,S}$.  It is therefore equivalent to isomorphism of their semisimplifications over $F$.  In particular it is equivalent to isomorphism when both representations are semisimple.

\begin{proof}
Suppose the trace functions differ.  Let $M$ be the $\mathcal O$-linear span of
\[
 \{(\rho_1(g),\rho_2(g)):g\in G_{\mathbb Q,S}\}
 \subseteq M_d(\mathcal O)\oplus M_d(\mathcal O).
\]
Products of spanning elements are again spanning elements, so $M$ is an $\mathcal O$-algebra with identity. It is a finite free $\mathcal O$-module of rank $r\le2d^2$, because it is a submodule of the finite free module on the right and $\mathcal O$ is a discrete valuation ring.  Reduction in this algebra gives a continuous finite quotient
\[
 \delta:G_{\mathbb Q,S}\longrightarrow H\subseteq(M/\varpi M)^\times.
\]
Continuity follows because $M$ is a finite module with its subspace topology and $\varpi M$ is open in $M$.  The quotient ring has $q^r$ elements; its zero element is not a unit. Hence $|H|\le q^r-1\le q^{2d^2}-1$.

Define the $\mathcal O$-linear functional
\[
 t:M\longrightarrow\mathcal O,\qquad t(A,B)=\operatorname{tr}A-\operatorname{tr}B.
\]
Its image is a nonzero ideal $\varpi^a\mathcal O$. Therefore $\lambda=\varpi^{-a}t$ is integral on $M$ and has nonzero reduction $\overline\lambda$ on $M/\varpi M$. Since the pairs $(\rho_1(g),\rho_2(g))$ span $M$, at least one has nonzero $\overline\lambda$-value. Conjugation by every such pair preserves $t$, hence also $\lambda$. Consequently
\[
 C=\{h\in H:\overline\lambda(h)\ne0\}
\]
is well defined, nonempty, and stable under conjugation. Every $g$ mapping to $C$ has unequal traces.

Let $L/\mathbb Q$ be the finite Galois extension defined by $\ker\delta$. It is unramified outside $S$, and $n_L=|H|\le q^{2d^2}-1$. Thus $\widetilde L=L(\sqrt m)$ has degree $\widetilde n\le N$ and is ramified only at primes dividing $m$. By the discriminant estimate of Serre, \emph{Quelques applications du th\'eor\`eme de densit\'e de Chebotarev}, Proposition~6 (also Mayle--Wang, Lemma~7),
\[
 \log D_{\widetilde L}
 \le(\widetilde n-1)\log\operatorname{rad}(D_{\widetilde L})
       +\widetilde n\log\widetilde n
 \le(N-1)\log m+N\log N.
\]
G01 supplies $p\nmid m$ of unequal traces with $p\le B(\widetilde L)\le\mathcal B(d,q,S)$. This proves the contrapositive of the finite-test assertion. The converse is immediate.  In characteristic zero, the Brauer--Nesbitt character criterion identifies equal trace functions with isomorphic semisimplifications.

The algebra construction is the extension to a coefficient discrete valuation ring of Mayle--Wang, Proposition~12, rather than an assumption of a known congruence depth. Its use of the quotient $M/\varpi M$ is essential: it must not be replaced by entrywise reduction of the ambient pair of matrix algebras.
\end{proof}

The exact analytic functions used by this proof are $\zeta_{\widetilde L}$ and $\zeta_{\mathbb Q}$, where $L$ is the finite quotient just constructed.  Neither the $L$-function of a positive-weight motive underlying $\rho_i$ nor Artin holomorphy is required. The finiteness and characteristic-zero character theory are independent of GRH. An executable numerical test additionally requires a way to compute the indicated exact traces; the theorem by itself does not manufacture that input.

\paragraph{Finding G04 (the precise limit of trace separation).}
Frobenius traces cannot distinguish arbitrary extension classes, even for lisse sheaves on an open subscheme of $\operatorname{Spec}\mathbb Z$.

\begin{proof}
Fix an odd prime $\ell$ and let $c(g)=\log_\ell(\chi_\ell(g)^{\ell-1})$, where $\chi_\ell$ is the cyclotomic character. This is a nonzero continuous additive homomorphism $G_{\mathbb Q,\{\ell\}}\to\mathbb Q_\ell$: its argument lies in $1+\ell\mathbb Z_\ell$, where the $\ell$-adic logarithm is an injective homomorphism, and the cyclotomic image contains this group. The two representations
\[
 \rho(g)=\begin{pmatrix}1&c(g)\\0&1\end{pmatrix},
 \qquad \rho_0(g)=\begin{pmatrix}1&0\\0&1\end{pmatrix}
\]
have trace $2$ and determinant $1$ on every element, including every unramified Frobenius. They are not isomorphic, because one has nontrivial image and the other has trivial image. Their semisimplifications are both the sum of two trivial characters.
\end{proof}

\paragraph{Finding G05 (explicit prime bounds for two LM7 predicates).}
Let $S$ be a finite set of rational primes and put
$m=\operatorname{rad}(2\prod_{p\in S}p)$.  Under the stipulated
Dirichlet GRH there is a prime $q\nmid m$ with $q\equiv2$ or $5\pmod9$
and
\[
 q\le\bigl[72\log3+24\log(4m)+35\bigr]^2. \tag{G05.1}
\]
Under $\mathrm H_M$ there is a prime $q\nmid m$ with $q\equiv4$ or
$7\pmod9$ and $3^{(q-1)/3}\not\equiv1\pmod q$, and
\[
 q\le\bigl[456\log3+72\log(4m)+95\bigr]^2. \tag{G05.2}
\]
These are quantitative refinements of the arithmetic predicates in the
recovered LM7 A.25 and A.16 registers, respectively.  Their source file
\texttt{sections/A\_agent.tex} was not recovered, so this finding does
not verify those rows' separate pro-three group presentations.

\begin{proof}
Put $L_0=\mathbb Q(\zeta_9)$ and $\alpha=\sqrt[3]{3}$.  The extension
$L_0/\mathbb Q$ is abelian of degree six.  The order
$\mathbb Z[\zeta_9]$ has discriminant of absolute value $3^9$:
for $\Phi_9(X)=X^6+X^3+1$,
\[
 \Phi_9'(\zeta_9)=3\zeta_9^2(2\zeta_3+1),\qquad
 N_{\mathbb Q(\zeta_3)/\mathbb Q}(2\zeta_3+1)=3.
\]
Taking the norm from $L_0$, whose degree over $\mathbb Q(\zeta_3)$ is
three, gives $3^9$.  Thus $D_{L_0}\le3^9$.  The two automorphisms
whose residues are $2,5\pmod9$ form the nonempty Frobenius predicate
for (G05.1).

The polynomial $X^3-3$ is Eisenstein at three.  Its real cubic field
$\mathbb Q(\alpha)$ is not Galois, so it cannot be contained in the
abelian extension $L_0$.  No other root can lie in $L_0$, either,
because $\zeta_3\in L_0$.  A cubic polynomial without a root is
irreducible; hence $L=L_0(\alpha)$ has degree eighteen over
$\mathbb Q$.  It is Galois, since $\zeta_3\in L_0$ makes it a
splitting field.  The product basis of
$\mathbb Z[\zeta_9]\mathbb Z[\alpha]$ is a basis of an order in $L$.
The cubic order has discriminant $-3^5$, and the tensor trace matrix
therefore gives
\[
 D_L\le (3^9)^3(3^5)^6=3^{57}. \tag{G05.3}
\]
For each $a=4,7$ and $j=1,2$ there is an automorphism
\[
 \zeta_9\longmapsto\zeta_9^a,\qquad
 \alpha\longmapsto\zeta_3^j\alpha.
\]
These four automorphisms form a nonempty conjugation-stable subset
of $\operatorname{Gal}(L/\mathbb Q)$.  At an unramified prime of residue
$a\pmod9$, the field $\mathbb F_q$ contains $\mu_3$, and $X^3-3$
has a root if and only if Frobenius fixes all its roots.  Thus the
four-element subset describes precisely the noncube predicate in
(G05.2).  Ordinary Chebotarev already gives density $4/18=2/9$;
the first predicate has density $2/6=1/3$.

For either $T=L_0$ or $T=L$, set $n=[T:\mathbb Q]$ and
$\widetilde T=T(\sqrt m)$.  Its degree is at most $2n$, and
\[
 \log D_{\widetilde T}\le2\log D_T+n\log(4m). \tag{G05.4}
\]
Indeed $m$ is positive, even and squarefree, so the quadratic field
has discriminant $4m$.  If it is not contained in $T$, product integral
bases give an order of discriminant $D_T^2(4m)^n$ in the degree-$2n$
compositum, whose field discriminant divides that order discriminant.
If it is contained, $\widetilde T=T$ and the displayed inequality
holds directly.  This covers the possible failure of linear disjointness.

Apply G01 and substitute (G05.4) in G-CHEB.  For $T=L_0$ we have
$[\widetilde T:\mathbb Q]\le12$ and
$\log D_{\widetilde T}\le18\log3+6\log(4m)$, giving (G05.1).
This compositum is abelian over $\mathbb Q$, so the stipulated
Dirichlet GRH supplies its Dedekind-zeta GRH.  For $T=L$ the
corresponding bounds are $36$ and $114\log3+18\log(4m)$, giving
(G05.2) under $\mathrm H_M$.  The chosen primes are unramified in
the relevant fields and avoid every divisor of $m$.  Neither use
requires Artin holomorphy.
\end{proof}

\subsection{Known explicit elliptic-curve consequences}

\paragraph{Cited result G-LOMBARDO (an unconditional bound over number fields).}
Let $E/K$ be an elliptic curve over a number field with
$\operatorname{End}_{\overline K}(E)=\mathbb Z$, let $d_K=[K:\mathbb Q]$,
and let $h_F(E)$ be its stable Faltings height in the convention of the
cited paper.  Lombardo, \emph{Bounds for Serre's open image theorem for
elliptic curves over number fields}, Algebra Number Theory 9 (2015),
2347--2395, Corollary~9.3, p.~2393, proves
\[
 [\operatorname{GL}_2(\widehat{\mathbb Z}):\rho_{E,\infty}(G_K)]
 < \exp(10^{21483})\,d_K^{\gamma}
       \max\{1,h_F(E),\log d_K\}^{2\gamma},
 \qquad \gamma=2.4\cdot10^{10}. \tag{G-LOMBARDO}
\]
This published result is unconditional. It bounds the adelic image index;
it does not assert surjectivity or identify the exact image. We use the
stated corollary, whose isogeny input is recorded in its Theorem~2.1,
and do not substitute the improvement in Remark~9.4 that invokes the
then-unpublished input of Remarks~2.3 and~2.7. For the actual number
field $L$ in LM8 A.32, it applies directly with $E=F_a$ and $K=L$.

\paragraph{Cited results G-MW1 and G-MW2.}
Mayle--Wang, \emph{On the effective version of Serre's open image theorem}, Bull. Lond. Math. Soc. 56 (2024), 1399--1416, Theorems~1 and~2, give the following conditional bounds. For a non-CM elliptic curve $E/\mathbb Q$, its mod-$\ell$ image equals $\operatorname{GL}_2(\mathbb F_\ell)$ whenever
\[
 \ell>964\log\operatorname{rad}(2N_E)+5760.
 \tag{G-MW1}
\]
For non-CM elliptic curves $E_1,E_2/\mathbb Q$ which are not $\mathbb Q$-isogenous, some prime of good reduction for both has unequal Frobenius traces and satisfies
\[
 p\le\left(482\log\operatorname{rad}(2N_{E_1}N_{E_2})+2880\right)^2.
 \tag{G-MW2}
\]
These become conditional consequences of $\mathrm{H}_M$: the analytic input in their \S3 and the proof of Theorem~2 is G-CHEB for finite auxiliary fields, together with the squarefree quadratic compositum of Corollary~6. The open-image theorem itself and Faltings's isogeny theorem were already unconditional. These numerical bounds are not bounds for arbitrary abelian varieties, arbitrary number fields, or CM elliptic curves. In particular they do not prove Serre's uniformity conjecture.

\paragraph{Finding G06 (the actual LM8 A.32--A.33 elliptic factors).}
The factors used in these rows have the models
\[
 F_a/\mathbb Q:\quad y^2=x^3-3x+a,\qquad a=6,-6.
\]
Under $\mathrm H_M$, their mod-$\ell$ Galois image over $\mathbb Q$ is
$\operatorname{GL}_2(\mathbb F_\ell)$ whenever
\[
 \ell>964\log6+5760. \tag{G06.1}
\]
For the number field $L$ chosen in the recovered proof, the restricted
image satisfies
\[
 [\operatorname{GL}_2(\mathbb F_\ell):
       \overline\rho_{F_a,\ell}(G_L)]\le[L:\mathbb Q]
 \quad\text{for every prime satisfying (G06.1).} \tag{G06.2}
\]
Independently of $\mathrm H_M$, G-LOMBARDO bounds their adelic image
over $L$. Since the models descend to $\mathbb Q$, restriction also gives
\[
 [\operatorname{GL}_2(\widehat{\mathbb Z}):\rho_{F_a,\infty}(G_L)]
 <[L:\mathbb Q]\,\exp(10^{21483})
          \max\{1,h_F(F_a)\}^{4.8\cdot10^{10}}. \tag{G06.3}
\]

\begin{proof}
The models and their role are explicit in recovered LM8 A.32.
Their non-CM hypothesis is supplied by the ordinary-reduction
argument of LM8 A.14: for $F_6$, the traces at $5,11$ are $-1,3$,
giving the distinct imaginary quadratic Frobenius fields with
discriminants $-19,-35$. A geometric characteristic-zero CM field
would inject into both, which is impossible. The other sign is a
quadratic twist. These two point counts were rechecked directly;
no classification of CM $j$-invariants is used.

The discriminant of each displayed integral model is
$-16(4(-3)^3+27\cdot6^2)=-2^9 3^3$.
The model has good reduction outside $2,3$. Its discriminant valuations
$9$ and $3$ are smaller than $12$, so no integral Weierstrass change
can lower them: discriminant valuations of isomorphic models differ
by a multiple of $12$. Both are positive in a minimal model, so both
primes are bad and $\operatorname{rad}(2N_{F_a})=6$.
Substitute this value in G-MW1.

For any homomorphism $\rho$ and subgroup $G_L\le G_{\mathbb Q}$,
the map on coset sets is surjective, giving
$[\rho(G_{\mathbb Q}):\rho(G_L)]\le[L:\mathbb Q]$.
With the mod-$\ell$ surjectivity just obtained this proves (G06.2).
With G-LOMBARDO over $\mathbb Q$, it proves (G06.3).
Full image after restriction to $L$ is not asserted. The original
rows use only qualitative commutator Zariski density, which was
already unconditional; these formulas supply numerical refinements
for their actual elliptic factors.
\end{proof}

\paragraph{Grothendieck framework and locators.}
For $X=\operatorname{Spec}\mathbb Z[1/\prod_{p\in S}p]$, finite quotients of $G_{\mathbb Q,S}$ are connected finite \'etale covers of $X$, and continuous integral adic representations are inverse systems of finite locally constant coefficient sheaves on $X$. The precise Galois-category references are SGA~1, Expos\'e~V, Th\'eor\`eme~4.1 and \S7, together with Propositions~8.1--8.2 (the generic-point description over a normal base). G01--G03 locate a closed point of this arithmetic scheme where a prescribed finite monodromy invariant is visible.  The finite quotient is a weight-zero Artin realization.  It is not a new purity statement for the original positive-weight sheaf.  This section supplies no theorem about crystalline realizations, divided powers, or general motivic $L$-functions; those structures are not used in these proofs.  Cote~149 scans~33--36 concern Kummer point separation and prime-to-characteristic geometric fundamental groups, not the numerical constants of these number-field prime bounds.

\paragraph{Remaining row-level gate.}
For further applications to programme rows, a numerical bound requires the finite quotient or pair of lisse sheaves, their coefficient ring, rank, ramification support, and the nonempty conjugacy-invariant test set. An open-image conclusion also requires the relevant geometric endomorphism and subgroup hypotheses. G06 verifies these model and non-CM inputs for the actual LM8 factors; G05 verifies the displayed arithmetic predicates while leaving their absent group-presentation proofs unverified. These distinctions are recorded in the row ledger.
